% ===== 原文 pp.28–32（PDF p.38–42） =====
% [原文 p.28]
\begin{align}
  \mathcal{D}\phi^I &= \dd\phi^I+A^aV^I_a \notag\\
  \mathcal{D}\zeta^I &= \dd\zeta^I+\Gamma^I_{JK}\mathcal{D}\phi^J\zeta^K+A^a\zeta^J\partial_JV^I_a \tag{3.5}
\end{align}
其中我们重新省略邻域标号 $(\alpha,\gamma)$，并理解为量 $\phi^I$ 与 $\zeta^I$ 只在 $M$ 上局部有定义。其它物质场的规范协变导数类似地定义。

我们还必须处理规范化的另一个潜在障碍。如上一节所写，4d A-模型的变分满足 $\delta_Q^2=0$；然而，若要成功地把该理论与一份 Donaldson-Witten 理论耦合，我们必须修改变分，使其满足 $\delta_Q^2=\delta_\sigma$，即参数为 $\sigma$ 的无穷小规范变换。这可以容易地通过在 $\zeta$、$\tilde\psi$、$\tilde H$ 的变分右端加入适当的依赖于 $V^I$ 的项来实现。加入这些项后，带规范场 4d A-模型的物质扇区的 BRST 变分为
\begin{align}
  \delta_Q\phi^I &= \zeta^I \notag\\
  \delta_Q\zeta^I &= \sigma^aV^I_a \notag\\
  \delta_QB^I &= \tilde\psi^I-\Gamma^I_{JK}\zeta^JB^K \notag\\
  \delta_Q\tilde\psi^I &= \frac12R^I{}_{JKL}B^J\zeta^K\zeta^L-\Gamma^I_{JK}\zeta^J\tilde\psi^K+\sigma^aB^J\nabla_JV^I_a \tag{3.6}\\
  \delta_Q\tilde\chi^I &= \tilde H^I-\Gamma^I_{JK}\zeta^J\tilde\chi^K \notag\\
  \delta_Q\tilde H^I &= \frac12R^I{}_{JKL}\tilde\chi^J\zeta^K\zeta^L-\Gamma^I_{JK}\zeta^J\tilde H^K+\sigma^a\tilde\chi^J\nabla_JV^I_a \notag
\end{align}
其中
\[
  \nabla_JV^I_a=\partial_JV^I_a+\Gamma^I_{JK}V^K_a.
\]

带规范场 4d A-模型的作用量 $S$ 取为如下各项之和：
\[
  S=\int_M\dd^4x\sqrt{h}\left(\mathcal{L}_1+\mathcal{L}_2+\mathcal{L}_3+\mathcal{L}_4+\mathcal{L}_5\right)-\frac{\ii\tau}{4\pi}\int_M\Tr F\wedge F,
\]

% [Tutor Note: p.28 “加入 $V^I$ 依赖的项以实现 $\delta_Q^2=\delta_\sigma$”可逐场验证：$\delta_Q^2\phi^I=\delta_Q\zeta^I=\sigma^aV^I_a=\delta_\sigma\phi^I$，$\delta_Q^2\zeta^I=\sigma^a\delta_QV^I_a=\sigma^a\zeta^J\partial_JV^I_a=\delta_\sigma\zeta^I$（与 (3.4) 一致）；而 $\delta_Q^2B^I$ 若不加上 $\sigma^aB^J\nabla_JV^I_a$ 则无法给出 $\delta_\sigma B^I=\sigma^aB^J\nabla_JV^I_a$（因 $\delta_\sigma$ 作用于张量场必须用协变导数 $\nabla$ 而不是偏导 $\partial$）。]

% [原文 p.29]
其中
\begin{align}
  \mathcal{L}_1 &= \frac{1}{e^2}\delta_Q\Tr\left\{\chi_{\mu\nu}\left(H^{\mu\nu}-2F^{+\mu\nu}\right)+2\bar\sigma D_\mu\psi^\mu\right\} \notag\\
  \mathcal{L}_2 &= \frac{1}{e^2}\delta_Q\left\{g_{IJ}\tilde\chi^I_\mu\left(\tilde H^{J\mu}-2\cD^\mu\phi^J-2\cD_\nu B^{J\nu\mu}\right)\right\} \notag\\
  \mathcal{L}_3 &= \frac{1}{e^2}\delta_Q\left\{-g_{IJ}\chi^a_{\mu\nu}\left(V^I_aB^{J\mu\nu}+\frac12(\nabla_KV^I_a)B^{J\mu\tau}B^{K\nu}_{\ \ \ \tau}\right)\right\} \tag{3.7}\\
  \mathcal{L}_4 &= \frac{1}{e^2}\delta_Q\left\{g_{IJ}\bar\sigma^a\left(2V^I_a\zeta^J+\frac12(\nabla_KV^I_a)\tilde\psi^J_{\mu\nu}B^{K\mu\nu}\right)\right\} \notag\\
  \mathcal{L}_5 &= \frac{1}{e^2}\delta_Q\Tr\left\{\eta[\sigma,\bar\sigma]\right\}. \notag
\end{align}

我们得到这个作用量的方式如下：首先，纳入 Donaldson-Witten 拓扑规范理论的作用量 $S_1$（含 $\theta$ 项）以及物质扇区的规范协变化作用量 $S_2$。由于我们的目标是与 Vafa-Witten 理论建立联系，我们选择纳入额外的“非极小”项 $S_3$、$S_4$、$S_5$，它们在靶空间与 $G$-作用取特定选择时约化为 Vafa-Witten 作用量中出现的项（下一节将证明）。特别地，最后一项 $S_5$ 是 BRST-正合的势项，它不影响辅助场的运动方程，也无助于局域化；它被纳入 Vafa-Witten 作用量，纯粹是为了使其在形式上等价于 $\cN=4$ SYM 作用量的扭转。

带规范场 4d A-模型的局域化方程为
\begin{gather*}
  F^{+a}_{\mu\nu}+\frac12g_{IJ}\kappa^{ab}V^I_bB^J_{\mu\nu}
    +\tfrac14g_{IJ}\kappa^{ab}(\nabla_KV^I_b)B^{J\tau}_{\ \ \ \mu}B^K_{\tau\nu}=0,\\
  D^\mu\phi^I+D_\nu B^{I\nu\mu}=0,
\end{gather*}
此外还有如下涉及复 0-形式 $\sigma$ 的方程：
\begin{gather*}
  D_\mu\sigma=0,\qquad [\sigma,\bar\sigma]=0,\\
  \sigma^aV^I_a=0,\qquad \sigma^aB^J_{\mu\nu}\nabla_JV^I_a=0.
\end{gather*}

% [Tutor Note: p.29 (3.7) 的 $\mathcal{L}_3$ 与 Vafa-Witten 形式之间的联系可这样逐项核对：代入 $X=\fg$、$g_{IJ}=-\kappa_{IJ}$、$V^I_a=f^I{}_{ab}\phi^b$（故 $\nabla_KV^I_a=f^I{}_{aK}$，靶空间平坦），则 $g_{IJ}\chi^a_{\mu\nu}V^I_aB^{J\mu\nu}=\chi^a_{\mu\nu}f^I{}_{aK}\phi^K B^{I\mu\nu}$ 型的缩并正是 $\Tr(\chi[B^{\mu\nu},\phi])$ 中的对易子项，而 $\frac12g_{IJ}\chi^a_{\mu\nu}f^I{}_{aK}B^{J\mu\tau}B^{K\nu}_{\ \ \ \tau}$ 则给出 $\frac12\Tr(\chi[B^{\mu\tau},B_\tau^{\ \nu}])$。整体符号取决于 $\Tr(XY)$ 与 $\kappa$ 的约定，见 p.32 的括注。]

\subsection{局部可观测量}

带规范场 4d A-模型的局部可观测量是 0-形式场 $\phi^I$、$\zeta^I$ 与 $\sigma$ 的规范不变且 BRST 不变的泛函。（检视各变分可以看出，BRST 上同调类中存在不含 0-形式场 $\bar\sigma$ 与 $\eta$ 的代表。）

% [原文 p.30]
给定这样的泛函 $\cO$，可以写出 $S(\fg^*)\otimes\Omega^\bullet(X)$ 中相应的\emph{等变形式}，其中 $S(\fg^*)\approx\C[\fg]$ 是 $\fg^*$ 上的对称代数，等价地也就是 $\fg$ 上的多项式代数。（关于等变形式的更多内容见 §C.2。）例如，我们映射
\[
  \cO=\frac{1}{k!}f_{I_1\ldots I_k}(\Phi)\,P(\sigma)\,\zeta^{I_1}\cdots\zeta^{I_k}
  \quad\longleftrightarrow\quad
  W_{\cO}=P(\sigma)\otimes\frac{1}{k!}f_{I_1\ldots I_k}(\Phi)\,\dd\phi^{I_1}\wedge\cdots\wedge\dd\phi^{I_k},
\]
其中 $P(\sigma)$ 是 $\sigma^a$ 的齐次多项式。若 $P$ 的次数为 $j$，则 $\cO$ 的鬼数为 $k+2j$。规范不变性把 $W_{\cO}$ 限制在 $G$-不变子代数
\[
  \left(S(\fg^*)\otimes\Omega^\bullet(X)\right)^G
\]
中。可以验证，带规范场 4d A-模型的 BRST 变分对应于作用等变 Cartan 微分
\[
  \dd_C=1\otimes\dd+\sigma^a\otimes\iota_a,
\]
其中 $\iota_a$ 表示用向量场 $V_a=V^I_a\dfrac{\partial}{\partial\phi^I}$ 缩并 $X$ 上的形式。因此，局部可观测量的 BRST 上同调对应于 $\Z$-分次复形
\[
  \left(\left(S(\fg^*)\otimes\Omega^\bullet(X)\right)^G,\ \dd_C\right)
\]
的上同调；由 Cartan 的一个定理 [16]，这恰好就是等变上同调 $H^\bullet_G(X)$。$\Z$-分次通过给 $\sigma^a$ 指派上同调次数 2 来实现，这与我们的鬼数指派一致。

% [Tutor Note: p.30 Cartan 微分 $\dd_C=1\otimes\dd+\sigma^a\otimes\iota_a$ 的幂零性需要 $G$-不变性条件：$\dd_C^2=\sigma^a(\dd\iota_a+\iota_a\dd)=\sigma^a\mathcal{L}_{V_a}$，而 $\mathcal{L}_{V_a}$ 作用在 $(S(\fg^*)\otimes\Omega^\bullet(X))^G$ 上为零（等价于 $1\otimes\dd$ 与 $\sigma^a\otimes\iota_a$ 的交叉项恰为李导数）。核对方法：$\sigma^a$ 与 $\iota_a$ 的分次交换性 + $[\iota_a,\iota_b]=0$ 使 $\sigma^a\sigma^b\iota_a\iota_b=0$。]

\section{Vafa-Witten 理论作为带规范场的 4d A-模型}

我们证明：对靶空间与 $G$-作用的特定选择，上一节的带规范场 4d A-模型等价于四维流形 $M$ 上规范群为 $G$ 的 Vafa-Witten 理论。令 $X$ 为向量空间 $\fg$，并取作用 $l$ 为 $G$ 在其李代数上的 Ad 作用。基 $T_a$ 给出 $\fg$ 的一组整体坐标 $\phi^a$；这进而消除了切指标 $I,J,\ldots$ 与李代数指标 $a,b,\ldots$ 之间的区别。$\fg$ 上的黎曼度量取为
\[
  g_{ab}=-\kappa_{ab}
\]
（回顾 $\kappa$ 是负定的）。由于 $X$ 是平坦向量空间，Christoffel 符号 $\Gamma^a_{bc}$ 恒为零，从而简化上一节变分与协变导数的表达式。对固定的 $\phi=\phi^aT_a\in\fg$，微分 $\dd l(-,\phi)$ 是 $\fg$ 在其自身上的 $\mathit{ad}$ 表示。相应的元素 $V\in T\fg\otimes\fg^*\simeq\fg\otimes\fg^*$ 的分量为
\[
  V^c_a=f^c{}_{ab}\phi^b.
\]

% [原文 p.31]
用矩阵语言，无穷小规范变换 (3.4) 约化为通常形式
\[
  \delta_\epsilon\phi=[\epsilon,\phi],\qquad \delta_\epsilon\zeta=[\epsilon,\zeta]
\]
（对取值于 $\mathit{ad}\,P$ 的场）。

规范扇区的变分 (3.2) 与物质扇区的变分 (3.6) 合起来约化为如下形式：
\begin{align}
  \delta_QA &= \psi, & \delta_Q\psi &= -\dd_A\sigma, \notag\\
  \delta_Q\sigma &= 0, \notag\\
  \delta_Q\bar\sigma &= \eta, & \delta_Q\eta &= [\sigma,\bar\sigma], \notag\\
  \delta_Q\chi &= H, & \delta_QH &= [\sigma,\chi], \tag{3.8}\\
  \delta_Q\phi &= \zeta, & \delta_Q\zeta &= [\sigma,\phi], \notag\\
  \delta_QB &= \tilde\psi, & \delta_Q\tilde\psi &= [\sigma,B] \notag\\
  \delta_Q\tilde\chi &= \tilde H, & \delta_Q\tilde H &= [\sigma,\tilde\chi]. \notag
\end{align}
确实，这些正是 [3] 与 [9] 中研究的 Vafa-Witten 拓扑规范理论离壳表述的场含量与变分。此外，作用量 (3.7) 约化为如下形式：
\[
  S=\int_M\dd^4x\sqrt{h}\left(\mathcal{L}_1+\mathcal{L}_2+\mathcal{L}_3+\mathcal{L}_4+\mathcal{L}_5\right)-\frac{\ii\tau}{4\pi}\int_M\Tr F\wedge F,
\]
其中
\begin{align}
  \mathcal{L}_1 &= \frac{1}{e^2}\delta_Q\Tr\left\{\chi_{\mu\nu}\left(H^{\mu\nu}-2F^{+\mu\nu}\right)+2\bar\sigma D_\mu\psi^\mu\right\} \notag\\
  \mathcal{L}_2 &= \frac{1}{e^2}\delta_Q\Tr\left\{\tilde\chi_\mu\left(\tilde H^\mu-2D^\mu\phi-2D_\nu B^{\nu\mu}\right)\right\} \notag\\
  \mathcal{L}_3 &= \frac{1}{e^2}\delta_Q\Tr\left\{-\chi_{\mu\nu}\left([B^{\mu\nu},\phi]+\frac12[B^{\mu\tau},B_\tau^{\ \nu}]\right)\right\} \tag{3.9}\\
  \mathcal{L}_4 &= \frac{1}{e^2}\delta_Q\Tr\left\{-\bar\sigma\left(\frac12[\tilde\psi_{\mu\nu},B^{\mu\nu}]+2[\zeta,\phi]\right)\right\} \notag\\
  \mathcal{L}_5 &= \frac{1}{e^2}\delta_Q\Tr\left\{\eta[\sigma,\bar\sigma]\right\} \notag
\end{align}
当 $M=\R^4$ 时，它与 $\cN=4$ SYM 理论的 Vafa-Witten 扭转的作用量 [9] 等价。

带规范场 4d A-模型的局域化方程约化为
\begin{align}
  F^+_{\mu\nu}+\frac12[B_{\mu\nu},\phi]+\tfrac14[B_\mu^{\ \tau},B_{\tau\nu}] &= 0 \tag{3.10}\\
  D^\mu\phi+D_\nu B^{\nu\mu} &= 0 \notag
\end{align}

% [原文 p.32]
此外还有方程
\begin{align}
  D_\mu\sigma &= 0 \notag\\
  [\sigma,\bar\sigma] &= [\sigma,\phi]=0 \tag{3.11}\\
  [\sigma,B_{\mu\nu}] &= 0. \notag
\end{align}
这些恰好等于 [3] 中的方程 (2.55) 与 (2.56)，后者定义了 Vafa-Witten 模问题。（为与 [3] 的记号对应，注意 Vafa 与 Witten 把 $\fg$-值场取为厄米而非反厄米的，并在他们的对易子中隐含了一个因子“$\ii$”。）[3] 中证明该作用量具有 $\phi\to-\phi$ 的对称性，因此上述方程组的每个解还额外满足
\[
  D_\mu\phi=[\phi,B_{\mu\nu}]=0.
\]

% [Tutor Note: p.32 末句的“由 $\phi\to-\phi$ 对称性推出 $D_\mu\phi=[\phi,B_{\mu\nu}]=0$”可直接核验对称性成立：在 $\phi\to-\phi$、$B_{\mu\nu}\to-B_{\mu\nu}$ 下，(3.10) 第一式三项分别变为 $F^+$、$+\frac12[B,\phi]$、$+\tfrac14[B,B]$（二次项变号两次），第二式整体变号，故解的集合在此变换下封闭。原书称由此可推出上述额外方程，完整论证见 [3]（需结合模空间的定义与规范固定）。]

% [Tutor Note: p.32 与 [3] 的记号对照是本节最易出错处：Vafa–Witten 用厄米（而非反厄米）$\fg$-值场，且在对易子中隐含因子 $\ii$，故 (3.10) 的 $\frac12$、$\tfrac14$ 系数与其文献形式看似不同。核对方法：把 $\phi\to\ii\phi$、$B\to\ii B$ 代入后，[3] 的方程 (2.55)–(2.56) 应逐项复现；若复现失败，多半是 $\kappa$ 与 $\Tr$ 的符号约定用错。]
